\begin{abstract}
Reward hacking is now observed in frontier reinforcement-learning runs, but it
is hard to study there: the true objective is rarely known exactly, and the
misspecification cannot be switched off. We present \framework{}, three
procedurally generated gridworld tasks in which the proxy reward differs from a
known true objective through one explicit, scalable corruption channel:
farmable reward shaping (\PF{}), tampering with a reward sensor (\ST{}), and a
bonus for detours that break a deadline the proxy omits (\GT{}). Across \res{n-runs-all} PPO
runs, hacking emerges in three different ways, from the late abandonment of a
mastered task to an exploit learned instead of the task. Solving
\res{optimal-layouts} layouts per task exactly shows that the proxy-optimal
policy hacks on most layouts from \res{opt-threshold-min}--\res{opt-threshold-max}
times the default channel strength, whereas PPO starts hacking only at several
times that in \PF{} and \GT{}; where PPO hacks in \PF{} and \ST{}, it gives up
more of the true objective than the optimum does. Monitors that see only
positions and actions reach held-out AUCs of \res{mon-behavior-auc-min}--\res{mon-behavior-auc-max},
but a baseline that knows only what kind of policy produced an episode reaches
\res{mon-behavior-identity-min}--\res{mon-behavior-identity-max}; within single
Proxy PPO policies, the monitors' AUC is
\res{mon-behavior-within-min}--\res{mon-behavior-within-max}, lowest where
hacking emerges late in training. Training agents against such a monitor
removes hacking in \PF{} and hides it in \GT{}, where agents still hack in \res{pen-GemTimeout-hack}\% of episodes while the
monitor's recall on hacks falls from \res{pen-GemTimeout-recall-before}\% to
\res{pen-GemTimeout-recall-after}\%. Every result cited in this paper is
regenerated from saved results by one command and checked in continuous
integration.
\end{abstract}
